\chapter{Surgery for Problems of the Small Toes}

\section*{What Is This Surgery?}
Surgery for problems of the small toes, often referred to as lesser toe surgery, is a group of orthopedic procedures designed to correct painful and debilitating deformities of the second, third, fourth, and fifth toes. These deformities include hammer toe, claw toe, and mallet toe, which involve abnormal flexion or extension at the metatarsophalangeal (MTP), proximal interphalangeal (PIP), or distal interphalangeal (DIP) joints. The primary goal of these surgical interventions is to restore proper anatomical alignment, alleviate pain, eliminate pressure points that cause corns and calluses, and improve the patient's ability to wear shoes comfortably and ambulate without discomfort. Procedures typically involve a combination of soft tissue releases, joint resections (arthroplasty), joint fusions (arthrodesis), and bone shortening or realignment osteotomies to achieve a straight, plantigrade toe. This surgical approach is crucial for patients whose quality of life is significantly impacted by chronic foot pain and functional limitations despite conservative management.

\section{Alternative Names}
\begin{itemize}
  \item Lesser Toe Deformity Correction
  \item Digital Arthroplasty
  \item Digital Arthrodesis
  \item Hammertoe Repair
  \item Clawtoe Correction
  \item Mallet Toe Reconstruction
\end{itemize}

\section*{Relevant Surgical Anatomy}
A thorough understanding of the intricate anatomy of the lesser toes and forefoot is paramount for successful surgical correction. The lesser toes consist of three phalanges (proximal, middle, and distal) for the second through fifth toes, articulating at the metatarsophalangeal (MTP), proximal interphalangeal (PIP), and distal interphalangeal (DIP) joints. The great toe (hallux) typically has only two phalanges. Key bony structures include the metatarsal heads, which articulate with the proximal phalanges, forming the MTP joints. The plantar plate, a fibrocartilaginous structure on the plantar aspect of the MTP joint, is critical for joint stability and preventing hyperextension. Collateral ligaments provide medial and lateral stability to all interphalangeal joints.

Soft tissue structures are equally important. The extensor digitorum longus and brevis tendons course dorsally, extending the toes, while the flexor digitorum longus and brevis tendons run plantarly, responsible for toe flexion. Imbalances in these tendon forces, often exacerbated by footwear or intrinsic foot deformities, contribute significantly to lesser toe deformities. The neurovascular supply to each toe is provided by dorsal and plantar digital nerves and arteries, which run along the sides of the phalanges. Surgical incisions must carefully avoid these bundles to prevent sensory deficits or vascular compromise. Lymphatic drainage generally follows the venous system. Surgical landmarks include the palpable joint lines of the MTP, PIP, and DIP joints, and the dorsal aspect of the toe where most incisions are made. The plantar plate is a critical structure to assess, particularly in cases of MTP joint instability.

\section{Surgical Principle \& Rationale}
The fundamental surgical principle for lesser toe deformities is to restore normal biomechanical alignment and balance the forces acting across the affected joints, thereby eliminating pain and preventing recurrence. The rationale stems from the understanding that these deformities arise from a combination of intrinsic muscle imbalance, extrinsic tendon contractures, and bony malalignment, often compounded by ill-fitting footwear. The operative concept involves addressing both the soft tissue contractures and the bony architecture. For hammertoe, characterized by hyperextension at the MTP joint and flexion at the PIP joint, the primary goal is to release the contracted extensor tendons and dorsal capsule at the MTP joint, and then to straighten the PIP joint. This straightening is typically achieved through either a PIP joint arthroplasty (resection of the proximal phalanx head) to create a flexible pseudoarthrosis, or an arthrodesis (fusion of the PIP joint) for rigid, stable correction.

The technical goals are multifaceted: to achieve a straight, stable, and plantigrade toe that rests flat on the ground; to eliminate painful pressure points, corns, and calluses; and to improve the patient's ability to wear conventional footwear comfortably. For claw toe, which involves hyperextension at the MTP joint and flexion at both the PIP and DIP joints, the approach is similar but may also require a flexor tenotomy or transfer, and potentially a DIP joint arthroplasty or arthrodesis. Mallet toe, characterized by isolated flexion at the DIP joint, is typically corrected with a DIP joint arthroplasty or arthrodesis. Internal fixation, often with Kirschner wires (K-wires) or small dedicated implants, is frequently employed to maintain the corrected position during the initial healing phase, ensuring the desired anatomical outcome.

\section{History \& Surgical Evolution}
Surgical correction of toe deformities has a rich history, evolving significantly since its early descriptions in the late 19th and early 20th centuries. Initial procedures were often rudimentary, focusing on simple resections of offending bony prominences or even partial amputations, which provided symptomatic relief but frequently resulted in unstable or flail toes. Early pioneers like Robert Jones in the UK described procedures for hammertoe correction in the early 1900s.

A significant milestone occurred in the mid-20th century when surgeons such as DuVries popularized the technique of proximal interphalangeal joint arthroplasty, involving the resection of the proximal phalanx head. This procedure became a cornerstone of hammertoe correction due to its relative simplicity and effectiveness. The introduction of Kirschner wires (K-wires) for temporary internal fixation further revolutionized these procedures, allowing for more stable and predictable fusions (arthrodesis) and resections, ensuring better maintenance of correction during healing. Further advancements in the late 20th and early 21st centuries focused on a deeper understanding of the complex biomechanics of the forefoot. This led to more sophisticated techniques, including balancing tendon forces, addressing MTP joint pathology (e.g., plantar plate repair, Weil osteotomy for metatarsal shortening), and the development of minimally invasive surgical approaches. The use of absorbable implants, specialized instrumentation, and a greater emphasis on preserving joint motion while achieving stable correction continue to refine the precision and outcomes of lesser toe surgery.

\section{Clinical Indications}
Surgery for problems of the small toes is indicated for specific clinical scenarios, typically after a thorough trial of conservative management has failed to provide adequate relief. These indications include:
\begin{itemize}
  \item Persistent pain in the affected toe or forefoot that significantly interferes with daily activities, ambulation, and the ability to wear conventional footwear.
  \item Development of painful corns or calluses on the dorsal aspect of the PIP joint, the tip of the toe, or the plantar aspect of the MTP joint due to chronic friction and pressure.
  \item Significant difficulty wearing conventional footwear due to the toe deformity, leading to limitations in lifestyle, social participation, and mobility.
  \item Progressive deformity that is becoming rigid or causing secondary problems such as MTP joint subluxation or dislocation, indicating structural instability.
  \item Recurrent ulceration or skin breakdown over bony prominences, particularly in patients with diabetes or neuropathy, where infection risk is high and wound healing is compromised.
  \item Functional impairment, such as instability during walking, inability to properly push off the ground, or altered gait mechanics.
  \item Cosmetic concerns when the deformity is severe and causes significant psychological distress, although this is rarely the sole indication for surgical intervention.
\end{itemize}

\section*{Clinical Positioning}
Surgery for problems of the small toes is generally considered a primary definitive treatment for symptomatic lesser toe deformities that have not responded to a comprehensive course of conservative care. Conservative measures, such as wider and deeper toe box shoes, custom orthotics, padding, strapping, activity modification, and anti-inflammatory medications, are typically attempted for several months (usually 3-6 months). If these non-surgical interventions fail to alleviate pain, improve function, or prevent skin breakdown, surgery becomes the primary option to correct the underlying structural and biomechanical problem.

In some specific contexts, lesser toe surgery may also serve as a secondary or salvage procedure. This can occur when a previous lesser toe surgery has failed, resulting in recurrence of the deformity, malunion, nonunion, persistent pain, or new iatrogenic deformities. In such cases, revision surgery is often more complex and requires careful planning. It can also be a secondary procedure when performed in conjunction with other forefoot surgeries, such as hallux valgus (bunion) correction or metatarsal osteotomies, to achieve overall forefoot balance, prevent transfer metatarsalgia (pain under adjacent metatarsal heads), and optimize the long-term outcome of the entire forefoot reconstruction.

\section{Surgical Technique \& Procedure}
Lesser toe surgery is typically performed on an outpatient basis under regional anesthesia (e.g., ankle block or popliteal block) with sedation, or general anesthesia. The patient is positioned supine, and a tourniquet is often applied to the ankle or thigh to provide a bloodless surgical field, enhancing visualization. The specific steps vary depending on the type and severity of the deformity and the joints involved, but generally follow a structured approach:

\begin{itemize}
  \item \textbf{Incision and Exposure:} A small dorsal longitudinal or transverse incision is made over the affected joint (e.g., PIP joint for hammertoe). For multiple toes, separate incisions are preferred. Meticulous care is taken to protect the underlying neurovascular bundles, which are typically retracted.
  \item \textbf{Soft Tissue Release:} For hammertoe and claw toe, the extensor tendon is often lengthened (extensor tenotomy) or released, and the dorsal joint capsule of the MTP joint may be incised (dorsal capsulotomy) to reduce hyperextension and decompress the joint.
  \item \textbf{Joint Arthroplasty (Resection):} For flexible or semi-rigid hammertoe and mallet toe, the head of the proximal phalanx (for PIP joint) or distal phalanx (for DIP joint) is resected using an osteotome or bone rongeurs. This creates a space that allows the toe to straighten and form a fibrous pseudoarthrosis.
  \item \textbf{Joint Arthrodesis (Fusion):} For rigid deformities, recurrent deformities, or when greater stability is desired, the articular surfaces of the PIP or DIP joint are resected, and the joint is fused in a corrected, slightly plantarflexed position. This provides a very stable, straight toe, albeit with permanent stiffness at that joint.
  \item \textbf{Bone Shortening/Osteotomy:} In cases of MTP joint subluxation, dislocation, or excessively long metatarsals contributing to the deformity, a shortening osteotomy of the metatarsal (e.g., Weil osteotomy) may be performed to decompress the MTP joint and restore proper length relationships.
  \item \textbf{Flexor Tendon Transfer:} For some claw toe deformities, the flexor digitorum longus tendon may be transferred dorsally to the proximal phalanx to help correct MTP joint hyperextension and provide a dynamic pull to keep the toe down.
  \item \textbf{Fixation:} After achieving correction, the toe is temporarily stabilized with a Kirschner wire (K-wire) inserted longitudinally through the toe and often into the metatarsal, or with small absorbable pins, staples, or screws. The K-wire is typically left protruding from the tip of the toe for easy removal later.
  \item \textbf{Closure:} The incision is meticulously closed in layers with absorbable sutures for deeper tissues and non-absorbable sutures or surgical glue for the skin. A sterile dressing is applied, and the toe is often taped or splinted to maintain alignment and provide initial support. No drains are typically used.
\end{itemize}

\section*{Preoperative Workup \& Preparation}
Thorough preoperative workup and preparation are crucial for optimizing outcomes and minimizing surgical risks. Patients undergo a detailed medical history and physical examination, including a comprehensive neurovascular assessment of the foot to evaluate circulation and sensation. Preoperative imaging typically includes weight-bearing X-rays of the foot (AP, lateral, and oblique views) to precisely assess the specific deformity, joint alignment, bone quality, and any associated forefoot pathologies. In some complex cases, an MRI may be ordered to evaluate soft tissue structures like the plantar plate or to rule out other pathologies.

Patients are strongly advised to stop smoking several weeks prior to surgery, as smoking significantly impairs wound healing, increases infection risk, and can lead to nonunion. Blood-thinning medications (e.g., aspirin, NSAIDs, warfarin, novel oral anticoagulants) must be stopped for a specified period before surgery, as directed by the surgeon and anesthesiologist, to minimize bleeding risk. Patients with diabetes require careful blood sugar control (HbA1c target <7.0\%) to reduce the risk of infection and poor wound healing. A preoperative visit with the anesthesia team will assess overall health and fitness for anesthesia, including cardiac and pulmonary evaluations. Patients are instructed to be NPO (nothing by mouth) for a specific duration (typically 6-8 hours) before surgery. Prophylactic antibiotics are often administered intravenously shortly before the incision to reduce the risk of surgical site infection. Finally, comprehensive informed consent is obtained, ensuring the patient fully understands the procedure, potential risks, expected benefits, and available alternatives.

\section*{Contraindications \& Precautions}
\textbf{Absolute Contraindications:}
\begin{itemize}
  \item Active infection in the surgical site (e.g., cellulitis, osteomyelitis) or systemic sepsis, which must be treated prior to elective surgery.
  \item Severe peripheral vascular disease (critical limb ischemia) that compromises blood supply to the foot, as this significantly impairs healing and increases the risk of tissue necrosis and amputation.
  \item Uncontrolled diabetes mellitus with severe neuropathy, poor glycemic control (e.g., HbA1c >8.0\%), or active Charcot neuroarthropathy, which predisposes to impaired wound healing, increased infection risk, and joint collapse.
  \item Non-ambulatory patient with no functional benefit expected from the procedure, as the goal is to improve mobility and comfort during weight-bearing.
  \item Unrealistic patient expectations regarding the outcome of the surgery, which can lead to dissatisfaction despite a technically successful procedure.
\end{itemize}

\textbf{Relative Contraindications \& Precautions:}
\begin{itemize}
  \item Poor nutritional status or significant immunosuppression (e.g., due to chronic illness, certain medications), which can delay wound healing and increase infection risk.
  \item Significant obesity (BMI >35-40), which may increase surgical complexity, anesthesia risks, and recovery time, and can place excessive stress on the healing foot.
  \item Smoking, which is a major risk factor for delayed wound healing, nonunion, and infection. Patients are strongly advised to cease smoking for at least 4-6 weeks preoperatively.
  \item Mild or asymptomatic toe deformities that do not cause pain, functional limitations, or skin breakdown, as conservative management is preferred in these cases.
  \item Certain rheumatologic conditions (e.g., severe rheumatoid arthritis) that may affect bone quality, soft tissue integrity, and healing capacity, requiring careful surgical planning and potentially modified fixation techniques.
  \item Previous failed toe surgery, requiring a more complex revision approach with potentially less predictable outcomes.
  \item Allergy to materials used for fixation (e.g., nickel in some implants), requiring careful selection of alternative hardware.
\end{itemize}

\section*{Complications \& Risks}
As with any surgical procedure, lesser toe surgery carries potential risks, which can be broadly categorized as general surgical risks and procedure-specific complications.

\begin{itemize}
  \item \textbf{General Surgical Risks:}
    \begin{itemize}
      \item \textbf{Bleeding:} Hematoma formation, though typically minor in foot surgery. Significant hemorrhage is rare.
      \item \textbf{Infection:} Superficial wound infection or deeper osteomyelitis, requiring prolonged antibiotic therapy or further surgical debridement.
      \item \textbf{Anesthesia Complications:} Nausea, vomiting, allergic reactions to medications, respiratory depression, cardiac events, or nerve damage from regional blocks.
      \item \textbf{Blood Clots:} Deep vein thrombosis (DVT) or pulmonary embolism (PE), though rare in foot surgery, especially with early ambulation.
    \end{itemize}
  \item \textbf{Procedure-Specific Risks:}
    \begin{itemize}
      \item \textbf{Stiffness of the Toe or Adjacent Joints:} Arthrofibrosis can limit range of motion, even after successful correction, particularly after arthrodesis.
      \item \textbf{Recurrence of the Deformity:} The toe may gradually return to its deformed position, especially if underlying biomechanical issues are not fully addressed or if fixation fails prematurely.
      \item \textbf{Malunion or Nonunion:} If an arthrodesis (fusion) is performed, the bones may heal in an incorrect position (malunion) or fail to fuse completely (nonunion), requiring revision surgery.
      \item \textbf{Floating Toe:} The corrected toe may not touch the ground during weight-bearing, leading to transfer pressure to adjacent metatarsal heads and potential new pain.
      \item \textbf{Nerve Injury:} Damage to the small dorsal or plantar digital nerves can result in persistent numbness, tingling, burning pain, or painful neuroma formation.
      \item \textbf{Vascular Compromise:} Although rare, injury to small digital blood vessels can impair blood supply to the toe, potentially leading to tissue loss or even partial toe amputation.
      \item \textbf{Hardware Complications:} K-wires or implants can migrate, break, cause irritation, become infected, or require premature removal due to discomfort.
      \item \textbf{Prolonged Swelling:} Swelling in the foot and toes can persist for several weeks to months postoperatively, delaying return to normal footwear.
      \item \textbf{Complex Regional Pain Syndrome (CRPS):} A rare but severe chronic pain condition affecting the limb, characterized by disproportionate pain, swelling, and skin changes.
      \item \textbf{Transfer Metatarsalgia:} Increased pressure and pain under an adjacent metatarsal head due to altered weight-bearing mechanics after correction of the primary deformity.
      \item \textbf{Unsatisfactory Cosmetic Result:} The toe may not appear perfectly straight, may have residual bulkiness, or may be shorter than desired.
    \end{itemize}
\end{itemize}

\section*{Postoperative Recovery Timeline}
Immediately after surgery, the patient is transferred to the post-anesthesia care unit (PACU) for close monitoring of vital signs, pain control, and neurovascular checks of the foot. The foot is typically elevated above heart level to minimize swelling and throbbing, and ice packs are applied. Most lesser toe surgeries are performed on an outpatient basis, allowing the patient to go home the same day. Upon discharge, patients are provided with prescriptions for oral pain medication, detailed instructions for wound care, and strict guidance on activity restrictions.

During the first 24-48 hours, strict elevation and regular ice application are crucial to manage swelling and pain. A special post-operative shoe or boot is typically worn to protect the surgical site and allow for protected weight-bearing, usually on the heel or flat-footed depending on the specific procedures performed. Stitches are usually removed at a follow-up visit around 10-14 days post-op. If K-wires were used for fixation, they are generally removed in the clinic between 3 to 6 weeks, once sufficient bone and soft tissue healing has occurred. Swelling and bruising are common and can persist for several weeks to months. Gradual return to normal activities and conventional shoe wear is guided by the surgeon, often taking 6-12 weeks for initial recovery, with residual swelling potentially lasting for several months. Physical therapy may be recommended to restore range of motion, strength, and improve gait mechanics.

\section*{Surgical Findings \& Pathology}
During lesser toe surgery, the surgeon meticulously documents intraoperative findings, which often include the specific type and rigidity of the deformity (e.g., flexible vs. rigid hammertoe), the degree of MTP joint hyperextension, the presence of plantar plate attenuation or tear, the extent of extensor tendon contracture, and any associated bony prominences or osteophytes. The condition of the articular cartilage at the resected joints is also noted. Any resected bone (e.g., proximal phalanx head, distal phalanx head) or soft tissue (e.g., hypertrophic capsule, scar tissue) is typically sent for pathological examination. The pathology report will confirm the nature of the tissue, often showing degenerative changes consistent with osteoarthritis, chronic inflammation, or fibrous tissue. It also serves to rule out any unexpected findings such as infection, tumor, or other specific disease processes, although these are rare in elective lesser toe surgery. These findings, combined with the successful correction achieved, form the basis for the operative report and contribute to the overall understanding of the patient's condition.

\section*{Expected Postoperative Course}
A normal and successful postoperative course following surgery for problems of the small toes is characterized by a progressive reduction in pain and swelling, leading to improved function and comfort. Initially, patients will experience moderate pain, managed with prescribed analgesics, which should gradually subside over the first few weeks. Swelling, though persistent for several months, will steadily decrease. The surgical incision should heal cleanly, typically within 2-3 weeks, without signs of infection. Once any temporary fixation (K-wires) is removed, the toe should maintain its corrected, straight, and plantigrade position, allowing it to lie flat on the ground during weight-bearing. Patients can expect a gradual return to protected weight-bearing, followed by increasing activity levels. The ultimate goal is the resolution of preoperative symptoms, including the absence of painful corns or calluses, and the ability to wear a wider variety of shoes comfortably, leading to enhanced mobility and an improved quality of life.

\section{Postoperative Care \& Follow-up}
Postoperative care and follow-up are essential to monitor healing, ensure optimal outcomes, and address any potential complications.
\begin{itemize}
  \item \textbf{Wound Checks:} Initial follow-up visits typically occur at 10-14 days post-surgery for wound inspection, dressing changes, and suture removal.
  \item \textbf{Pin Removal:} If Kirschner wires (K-wires) were used for fixation, they are usually removed in the clinic between 3 to 6 weeks post-surgery, once sufficient bone and soft tissue healing has occurred. This is typically a quick, minimally painful procedure.
  \item \textbf{Imaging:} Post-operative X-rays may be taken at various intervals (e.g., at 6 weeks and 3 months) to confirm proper alignment and assess bone healing, especially if an arthrodesis (fusion) was performed.
  \item \textbf{Physical Therapy/Rehabilitation:} Depending on the specific procedure and individual recovery, physical therapy may be recommended to improve range of motion, strength, and gait mechanics, particularly after MTP joint procedures or if stiffness persists.
  \item \textbf{Activity Resumption:} Patients are guided on a gradual return to normal activities, including exercise and sports, typically over several months, with full recovery often taking 3-6 months.
  \item \textbf{Footwear Advice:} Guidance on appropriate footwear (e.g., wider toe box shoes, avoiding high heels) is provided to prevent recurrence and maintain comfort.
\end{itemize}
Patients are instructed to contact their surgeon immediately if they experience increased pain, redness, swelling, warmth, fever, purulent drainage from the wound, or any changes in sensation or color of the toe, as these could indicate infection or other complications.

\section*{Epidemiology}
Lesser toe deformities are among the most common orthopedic conditions affecting the foot, with their prevalence increasing significantly with age. Hammertoe is the most frequently encountered deformity, often affecting the second toe, and is more common in women, largely attributed to the prolonged use of narrow, high-heeled footwear that compresses the forefoot. The incidence of these deformities is estimated to be as high as 20-30\% in the adult population, with symptomatic cases requiring intervention being a smaller but substantial subset. These conditions are frequently associated with other forefoot pathologies, such as hallux valgus (bunion), pes planus (flatfoot), and metatarsalgia, which contribute to altered biomechanics and muscle imbalances. While mortality directly related to elective lesser toe surgery is exceedingly rare, morbidity can include recurrence rates ranging from 10-20\%, infection rates of 1-5\%, and persistent stiffness or pain in a similar percentage of patients. The economic burden includes direct healthcare costs associated with surgical procedures and indirect costs related to lost productivity and reduced quality of life due to chronic pain and disability.

\section*{Outcomes \& Prognosis}
The overall outcomes and prognosis for surgery for problems of the small toes are generally favorable, with high rates of patient satisfaction. Success rates, defined by significant pain relief, improved function, and the ability to wear shoes comfortably, typically range from 80\% to 90\%. Patients often experience improved mobility, enhanced participation in daily activities, and a better overall quality of life. Functional outcomes are usually excellent, with most patients returning to their desired activity levels within 3-6 months, although residual swelling can persist longer. Recurrence of the deformity is the most common long-term complication, occurring in approximately 10-20\% of cases, particularly with flexible deformities or if the underlying biomechanical imbalances are not fully addressed. Prognostic factors influencing outcomes include the severity and rigidity of the initial deformity, the presence of associated forefoot pathologies (e.g., MTP joint instability, hallux valgus), patient compliance with post-operative instructions, and the absence of systemic conditions like severe diabetes or peripheral vascular disease. While fusion procedures (arthrodesis) tend to have lower recurrence rates due to rigid fixation, they result in a permanently stiff joint, which some patients may perceive as a functional limitation.

\section{References}
\begin{enumerate}
  \item MedlinePlus. Hammertoe. U.S. National Library of Medicine; 2024. Available from: \url{https://medlineplus.gov/hammertoe.html}
  \item Mann RA, Coughlin MJ, Saltzman CL, et al. Mann's Surgery of the Foot and Ankle. 9th ed. Elsevier; 2023.
  \item American Academy of Orthopaedic Surgeons (AAOS). Hammertoe. OrthoInfo; 2023. Available from: \url{https://orthoinfo.aaos.org/en/diseases--conditions/hammertoe/}
  \item Campbell WC, Canale ST, Beaty JH. Campbell's Operative Orthopaedics. 14th ed. Elsevier; 2021.
\end{enumerate}

\clearpage